\chapter{Thực nghiệm}
\label{ch:experiments}

Chương này (1) mô tả môi trường thực nghiệm và quy trình đánh giá, (2) báo
cáo kết quả thu thập dữ liệu, (3) báo cáo kết quả fine-tuning, (4) trình
bày đánh giá chính trên bộ task chính thức của bốn website WebArena, và
(5) phân tích kết quả.

% GHI CHÚ: chương này dựng khung trước và sẽ được bổ sung khi các thực
% nghiệm còn lại hoàn tất (các mốc [BỔ SUNG]). Số liệu hiện có lấy từ đợt
% đánh giá GitLab ngày 15/07/2026.

\section{Môi trường thực nghiệm}
\label{sec:exp-setup}
% TODO: các instance WebArena tự triển khai; vast.ai A100-80GB cho huấn
% luyện và serving; bảng cấu hình LLM theo vai trò; tham số run; đánh giá
% qua BrowserGym scored environments + bộ chấm chính thức, phát lại hành
% động trả lời cuối.
\textit{[Sẽ viết — xem dàn ý 07, mục 4.1.]}

\section{Kết quả thu thập dữ liệu}
\label{sec:exp-data}
% TODO: thống kê từng run (node, task đề xuất/giải được/bất khả thi,
% skill, operator, chuỗi); dataset cuối: 370 mẫu (352 train / 18 val);
% 1-2 ví dụ tiêu biểu về skill và semantic task.
\textit{[Sẽ viết — xem dàn ý 07, mục 4.2.]}

\section{Kết quả fine-tuning}
\label{sec:exp-sft}
% TODO: đường loss (train 0.32 / eval 0.24), bảng thời gian và chi phí.
\textit{[Sẽ viết — xem dàn ý 07, mục 4.3.]}

\section{Đánh giá chính trên WebArena}
\label{sec:exp-main}
% Đánh giá trên BỐN website: Shopping, Shopping Admin, GitLab, Reddit.
% Hiện có: GitLab — mô hình gốc 10/180 (5,56%) vs fine-tune 12/180 (6,67%).
% Trình bày là tín hiệu định hướng (mỗi task chạy 1 lần) cho tới khi có 4.5.
% [BỔ SUNG] kết quả Shopping / Shopping Admin / Reddit + bảng theo website.
\textit{[Sẽ viết — xem dàn ý 07, mục 4.4.]}

\section{Lặp lại thống kê}
\label{sec:exp-stats}
% [BỔ SUNG] chạy lặp >=3-5 lần (hoặc nhiều seed) mỗi mô hình -> mean +/- SD,
% 95% CI, kiểm định phi tham số. Ưu tiên số 1.
\textit{[Dự kiến — xem dàn ý 07, mục 4.5.]}

\section{Nghiên cứu ablation}
\label{sec:exp-ablation}
% [BỔ SUNG] (a) dữ liệu khám phá theo graph vs khám phá phẳng cùng ngân
% sách; (b) có/không judge; (c) có/không viết lại suy nghĩ; (d) phân tích
% liên-website từ kết quả 4 website.
\textit{[Dự kiến — xem dàn ý 07, mục 4.6.]}

\section{Phân tích}
\label{sec:exp-analysis}
% TODO: phân rã theo loại task, ví dụ thành công/thất bại định tính,
% phân loại lỗi (nối sang Chương 5).
\textit{[Sẽ viết — xem dàn ý 07, mục 4.7.]}
